\pg{342} 65. O finito é cognoscivel em \textbf{si mesmo}. --- O \emph{finito}, emquanto tal, importa duas coisas: a \emph{entidade} e a \emph{limitação}. Ora nós podemos muito bem conhecer, em si mesma, a \emph{entidade}, que é uma coisa absoluta, como podemos conhecer, na sua entidade, a \emph{limitação}, que é apenas a negação de ulterior perfeição. Logo o \emph{finito} é cognoscivel em si mesmo \footnote{\emph{Descartes}, como vimos, sustentou que a nossa intelligencia concebe primeiramente a idea do \emph{infinito}, e depois a do \emph{finito}; porque, na opinião d'elle, não poderiamos dizer que um ente é limitado, se não tivessemos já a idea do ente \emph{illimitado}. --- Mas é manifesto o seu erro. O \emph{finito} é a. \emph{entidade limitada}. Ora, como dizemos no texto, tanto a \emph{entidade}, quanto a\pgn{343}\emph{limitação}, podem ser percebidas em si mesmas, sem que primeiramente percebamos o \emph{infinito}; porque a \emph{entidade} e a \emph{limitação} nem importam infinidade, nem se referem ao infinito.}.

- D

